\section{核心：Newton--Schulz 迭代与 Muon}
\label{sec:ns-muon}

\subsection{Newton--Schulz 迭代}

设 $G \in \mathbb{R}^{m\times n}$（$m \ge n$，列满秩），考虑迭代
\begin{equation}
  X_0 = \frac{G}{\|G\|_2},\qquad X_{k+1} = \tfrac{1}{2} X_k (3 I - X_k^\top X_k).
  \label{eq:ns}
\end{equation}
这就是 Schulz (1933) 提出的 Newton 型矩阵函数迭代 \cite{schulz1933,higham2008}。在标量情形下退化为 $x_{k+1} = x_k(3 - x_k^2)/2$。

\begin{algorithm}[htbp]
  \caption{Newton--Schulz 极分解迭代}
  \label{alg:ns}
  \begin{algorithmic}[1]
    \Require 矩阵 $G\in\mathbb{R}^{m\times n}$（列满秩）
    \Ensure 正交因子近似 $X_k \approx U$
    \State $X_0 \gets G / \|G\|_2$
    \For{$k = 0,1,2,\ldots$}
      \State $X_{k+1} \gets \tfrac{1}{2}\, X_k \,(3I - X_k^\top X_k)$
    \EndFor
  \end{algorithmic}
\end{algorithm}

\begin{theorem}[局部二次收敛]
\label{thm:ns}
若 $X_0$ 满足 $\sigma_{\min}(X_0) > 0$ 且 $\sigma_{\max}(X_0) < \sqrt 3$，则 \eqref{eq:ns} 二次收敛到 $G$ 的极分解正交因子 $U$。令 $E_k = X_k^\top X_k - I$，则
\begin{equation}
  E_{k+1} = -\tfrac{1}{4} E_k^2 \left(3 I - E_k\right),
  \label{eq:ns-error}
\end{equation}
故 $\|E_{k+1}\|_F \le \tfrac{1}{4}\|E_k\|_F^2 (3 + \|E_k\|_F)$，即 $\|E_k\|_F$ 二次衰减。
\end{theorem}

\begin{proof}
设 $G_k = X_k^\top X_k$。由 \eqref{eq:ns}，$G_{k+1} = \tfrac{1}{4} (3 I - G_k) G_k (3 I - G_k)$。代入 $G_k = I + E_k$，展开化简得 $E_{k+1} = -\tfrac{1}{4} E_k^2 (3 I - E_k)$。取 Frobenius 范数，当 $\|E_k\|_F < 1$ 时 $\|E_{k+1}\|_F < \|E_k\|_F^2$，即二次收敛。
\end{proof}

\textbf{收敛盆的标量约简}：注意迭代 \eqref{eq:ns} 只通过 $X_k^\top X_k$ 作用于 $X_k$，因此它在奇异值上是\textbf{可对角化}的。设 $X_k = P_k \Sigma_k Q_k^\top$ 是 SVD，则 $X_k^\top X_k = Q_k \Sigma_k^2 Q_k^\top$，代入 \eqref{eq:ns} 得
\begin{equation}
  X_{k+1} = \tfrac{1}{2} P_k \Sigma_k\,(3I - \Sigma_k^2)\,Q_k^\top
          = P_k\,\varphi(\Sigma_k)\,Q_k^\top,
  \label{eq:ns-svd}
\end{equation}
其中 $\varphi$ 逐元素作用于奇异值。故 $X_k$ 始终保持初始奇异向量 $P_0,Q_0$ 不变，每个奇异值 $\sigma$ 独立地按标量映射
\begin{equation}
  \varphi(\sigma) = \tfrac{1}{2}\,\sigma\,(3 - \sigma^2)
  \label{eq:ns-scalar}
\end{equation}
演化。于是矩阵迭代收敛到正交因子 $U=P_0 Q_0^\top$（即所有奇异值 $\to 1$）当且仅当每个初始奇异值都落在标量映射 \eqref{eq:ns-scalar} 收敛到 $1$ 的吸引盆内。这把定理~\ref{thm:ns} 的收敛盆问题严格化为对 $\varphi$ 的一维分析。

\begin{proposition}[标量收敛盆 $(0,\sqrt 3)$]
\label{prop:basin}
设 $\sigma_0 \in (0,\sqrt 3)$，按 \eqref{eq:ns-scalar} 定义 $\sigma_{k+1}=\varphi(\sigma_k)$，则 $\sigma_k \in (0,\sqrt 3)$ 对所有 $k$ 成立，且 $\sigma_k \to 1$；收敛是二次的：$|\sigma_{k+1}-1| \le \tfrac{3}{2}\,|\sigma_k - 1|^2$ 当 $\sigma_k$ 足够接近 $1$。若 $\sigma_0 > \sqrt 3$，则 $\sigma_1 < 0$ 且 $|\sigma_k|\to\infty$。
\end{proposition}

\begin{proof}[证明（要点）]
在 $(0,\infty)$ 上 $\varphi$ 的不动点由 $\varphi(\sigma)=\sigma$ 即 $\sigma(1-\sigma^2)=0$ 给出，正实轴上唯一非零不动点为 $\sigma^*=1$；又 $\varphi'(\sigma)=\tfrac{3}{2}(1-\sigma^2)$，故 $\varphi'(1)=0$，即 $\sigma^*=1$ 是\textbf{超吸引}不动点（这正是二次收敛的来源）。

\emph{区间 $(0,1]$}：此处 $\varphi'(\sigma)=\tfrac{3}{2}(1-\sigma^2)\ge 0$，$\varphi$ 单调不减，$\varphi(0)=0$、$\varphi(1)=1$，故 $\varphi$ 把 $(0,1]$ 映入 $(0,1]$。又 $\varphi(\sigma)-\sigma=\tfrac{1}{2}\sigma(1-\sigma^2)>0$ 对 $\sigma\in(0,1)$ 成立，于是 $\{\sigma_k\}$ 单调递增且以 $1$ 为上界，必收敛；极限只能是不动点 $1$。

\emph{区间 $(1,\sqrt 3)$}：此处 $\varphi$ 单调递减，$\varphi(1)=1$、$\varphi(\sqrt3)=0$，故 $\varphi\bigl((1,\sqrt3)\bigr)=(0,1)$。因此任何 $\sigma_0\in(1,\sqrt3)$ 一步后落入 $(0,1)$，随后归结为上一情形，单调递增收敛到 $1$。综合两段，$(0,\sqrt3)$ 是 $1$ 的吸引盆且迭代始终保持在 $(0,\sqrt3)$ 内。

\emph{二次速率}：令 $e_k=\sigma_k-1$，由 $\varphi(\sigma)=1-\tfrac{3}{2}(\sigma-1)^2-\tfrac{1}{2}(\sigma-1)^3$ 直接展开得 $e_{k+1}=-\tfrac{3}{2}e_k^2-\tfrac{1}{2}e_k^3$，故 $|e_{k+1}|\le\tfrac{3}{2}|e_k|^2$（$|e_k|\le1$ 时），与正文递推 $E_{k+1}=-\tfrac{1}{4}E_k^2(3I-E_k)$ 在标量化后一致。

\emph{发散侧}：若 $\sigma_0>\sqrt3$ 则 $3-\sigma_0^2<0$，故 $\sigma_1=\tfrac12\sigma_0(3-\sigma_0^2)<0$；此后 $|\varphi(\sigma)|\sim\tfrac12|\sigma|^3$ 主导，$|\sigma_k|\to\infty$。$\square$
\end{proof}

\noindent 命题~\ref{prop:basin} 给出定理~\ref{thm:ns} 中收敛盆断言的完整证明；完整推导见附录~\ref{app:proofs}。

\textbf{实验 30}：在 100 个初值 $\sigma_0 \in [0.05, 2.5]$ 上跑 NS，落在 $[0.05, \sqrt 3)$ 内的初值全部收敛；落在 $(\sqrt 3, 2.5]$ 的 11 个初值全部发散，与命题~\ref{prop:basin} 的预言完全一致（图~\ref{fig:exp30}）。

\begin{figure}[htbp]
  \centering
  \IfFileExists{../figures/exp30_ns_basin.png}{%
    \includegraphics[width=0.85\linewidth]{exp30_ns_basin.png}%
  }{%
    \fbox{\parbox{0.8\linewidth}{\centering\vspace{2cm}图 exp30：运行实验后生成\\（\texttt{experiments/run\_all.py}）\vspace{2cm}}}%
  }
  \caption{Newton--Schulz 收敛盆：$\sigma_0$ 与 40 步后 $|\sigma^2-1|$}
  \label{fig:exp30}
\end{figure}

\subsection{五次多项式 Higham 加速}

实验 27 用 Higham (2008, eq.~8.20) 五次多项式
\[
  X_{k+1} = \tfrac{1}{8}\bigl(15 X_k - 10 X_k(X_k^\top X_k) + 3 X_k(X_k^\top X_k)^2\bigr),
\]
可验证 $\phi(1)=1$、$\phi'(1)=\phi''(1)=0$，5 步达机器精度，标准三阶 NS 要 11--15 步 \cite{higham2008,anonymous2025ns}。

\subsection{Muon 优化器：谱范数 trust-region}

Muon 的更新规则 \cite{jordan2024muon}：

\begin{algorithm}[htbp]
  \caption{Muon 优化器（矩阵参数）}
  \label{alg:muon}
  \begin{algorithmic}[1]
    \Require 矩阵参数 $W$，损失 $f(W)$，步长 $\eta$，动量系数 $\beta\in[0,1)$，Newton--Schulz 步数 $T_{\mathrm{ns}}$
    \Ensure 更新后的 $W_{k+1}$
    \State 初始化动量缓冲 $B_0 \gets 0$
    \State $G_k \gets \nabla f(W_k)$；\quad $B_{k+1} \gets \beta B_k + G_k$
    \State $M_k \gets G_k + \beta B_{k+1}$ \Comment{Nesterov 动量梯度；$\beta=0$ 时 $M_k=G_k$}
    \State $\widetilde M_k \gets \textsc{NewtonSchulz}(M_k,\, T_{\mathrm{ns}})$ \Comment{算法~\ref{alg:ns}，迭代 $T_{\mathrm{ns}}$ 步}
    \State $W_{k+1} \gets W_k - \eta\, \widetilde M_k$
  \end{algorithmic}
\end{algorithm}

\noindent 本文实验中 Newton--Schulz 步数取 $T_{\mathrm{ns}}=8$（矩阵二次/谱范数实验）或 $5$（真实数据训练，见 \S\ref{sec:experiments} 实验 31）；动量系数在受控对照实验（实验 12、28）中取 $\beta=0$，即直接对原始梯度做正交化以隔离正交化方向的几何效应，在真实训练实验中取 $\beta=0.9$。

\begin{theorem}[极分解的 trust-region 最优性]
\label{thm:muon-tr}
对任意 $G \in \mathbb{R}^{m\times n}$（$\mathrm{rank}(G) = r > 0$），
\[
  -t\,U V^\top \in \arg\min_{\substack{\Delta \in \mathbb{R}^{m\times n}\\ \|\Delta\|_2 \le t}} \langle G, \Delta\rangle_F,
\]
其中 $G = U \Sigma V^\top$ 是紧 SVD \cite{kovalev2025}。
\end{theorem}

\begin{proof}
设 $\Delta = \widetilde U \widetilde\Sigma \widetilde V^\top$，$\|\Delta\|_2 = \widetilde\sigma_{\max} \le t$。由 von Neumann 迹不等式，
\[
  |\mathrm{tr}(\Delta^\top G)| \le \sum_{i=1}^{r} \sigma_i(\Delta) \sigma_i(G) \le t \sum_i \sigma_i(G) = t\,\|G\|_*.
\]
等号由 $\Delta = -t U V^\top$ 取到。
\end{proof}

\textbf{含义}：Muon 在谱范数 trust-region 上做线性逼近最速下降。这就是为什么 Muon 与 Frobenius GD 在 Frobenius 损失上不必一致。

\subsection{Muon 在匹配谱范数几何的目标上下降}

为展示 Muon 的优势，实验 12 和 28 设计了\textbf{谱范数恢复目标}
\[
  \min_{W \in \mathbb{R}^{m\times n}} f(W) = \tfrac{1}{2} \|W - W^*\|_\sigma^2,
\]
其中 $\|\cdot\|_\sigma$ 是谱范数。

\paragraph{实验设置与可复现细节。}\quad
谱范数恢复实验中 $W^*$ 由 $W^*=U\Sigma V^\top$ 生成：$U\in\mathbb{R}^{m\times r}$、$V\in\mathbb{R}^{n\times r}$（$r=\min(m,n)$）取标准正态随机矩阵的 QR 正交因子（Haar 分布），奇异值 $\sigma_i\stackrel{\text{iid}}{\sim}\mathrm{Uniform}[0.5,1.5]$。所有方法均从 $W_0=0$ 出发。实验 12 取维度 $m\times n=16\times 8$、迭代 $80$ 步，Muon 步长 $\eta=0.3$、$T_{\mathrm{ns}}=8$、$\beta=0$，GD 步长 $0.5$，Adam 步长 $0.3$（默认 $(\beta_1,\beta_2)=(0.9,0.999)$）。实验 12 右图的 Frobenius 对照取矩阵二次 $\tfrac12\|AW-B\|_F^2$，$A=Q\,\mathrm{diag}(\mathrm{geomspace}(1,100))\,Q^\top$（$\kappa(A)=100$，$Q$ 为 Haar 正交阵），$W^*$ 元素 iid $\mathcal N(0,1)$、$B=AW^*$，维度 $16\times 8$、$400$ 步，Muon 取 $\eta=10^{-3}/L$、$T_{\mathrm{ns}}=8$。实验 28 取 $m\times n=12\times 8$、$60$ 步、$5$ 个随机种子，超参与实验 12 谱范数组一致。全部随机数由固定种子（主实验 \texttt{seed}=42，多种子实验 \texttt{seed}$\in\{0,\dots,4\}$）控制。实验 12 左图（80 步、$W_0=0$）：Muon-NS 从 1.077 降到 \textbf{0.037}；GD 终值 $2.4\times 10^{-32}$、Adam 终值 $1.1\times 10^{-4}$。注意 $W_0=0$ 时 $\nabla f(0)=-W^*$，GD 沿 $W^*$ 直线下降，故左图中 GD 曲线往往最低——这是初值巧合，详见 \S\ref{sec:experiments} 对图~\ref{fig:exp28} 的读图说明。实验 28 在 5 个种子上 Muon 最终损失 0.028$\sim$0.037 高度一致。

\textbf{与 Frobenius 矩阵二次的对照}（exp12 右图）：在 $\min_W \tfrac{1}{2}\|AW - B\|_F^2$ 上 Frobenius GD 降到 $f - f^* = 51$，Muon 反而高达 63266——这才是 Muon 与 GD 几何错配的典型表现，也是定理~\ref{thm:muon-tr} 的直接含义。

\begin{figure}[H]
  \centering
  \IfFileExists{../figures/exp12_matrix_muon.png}{%
    \includegraphics[width=0.9\linewidth]{exp12_matrix_muon.png}%
  }{%
    \fbox{\parbox{0.8\linewidth}{\centering\vspace{2cm}图 exp12\vspace{2cm}}}%
  }
  \caption{Muon 在谱范数损失 vs Frobenius 矩阵二次上的表现（左：$W_0=0$ 时 GD 因 $\nabla f(0)=-W^*$ 往往最低）}
  \label{fig:exp12}
\end{figure}

\subsection{如何理解 Muon 与 GD 的对比}
\label{sec:muon-vs-gd}

读者容易得出``Muon 各种角度不如 GD''的印象，这源于\textbf{用错了比较标尺}：Muon 不是 Frobenius 意义下的 GD 替代品，而是换了一套预条件几何 $P_k$（矩阵正交化 + 谱范数 trust-region）。二者在统一模板 $x_{k+1}=x_k-P_k g_k$ 下\textbf{可比的是机制与理论}，不宜在同一损失、同一初值下简单比``谁更小、谁更快''。

\begin{table}[H]
  \centering
  \caption{Muon 与 GD 在不同实验中的角色对照}
  \label{tab:muon-gd-roles}
  \small
  \begin{tabularx}{\textwidth}{@{}lXX@{}}
    \toprule
    实验场景 & GD 的表现 & Muon 应如何解读 \\
    \midrule
    向量强凸二次（exp1） & 基准方法，收敛慢但可靠 & 不适用；Muon 面向矩阵参数 \\
    Frobenius 矩阵二次（exp12 右） & 正常下降至 $f-f^*=51$ & \textbf{发散}至 $63266$：说明 Muon 不能替代 Frobenius GD \\
    谱范数恢复（exp12 左 / exp28） & $W_0{=}0$ 时常因 $\nabla f(0){=}{-}W^*$ 直达 $10^{-32}$ & 稳定降至 $0.03$ 平台：验证 trust-region \textbf{方向几何}匹配，非全局精度 \\
    \bottomrule
  \end{tabularx}
\end{table}

\noindent\textbf{Muon 的价值不在``全面打败 GD''，而在：}
\begin{enumerate}[nosep]
  \item \textbf{理论闭环}：核心子程序 Newton--Schulz 是课内矩阵函数迭代；定理~\ref{thm:muon-tr} 给出每步在谱范数 trust-region 下的最优方向（von Neumann 迹不等式取等）。
  \item \textbf{几何分工}：GD 优化 Frobenius 线性子问题；Muon 优化谱范数 trust-region 线性子问题——目标不同，胜负不可直接横比。
  \item \textbf{工程动机} \cite{jordan2024muon,kovalev2025}：LLM 训练中矩阵型隐层权重需要在谱范数意义下控制更新幅度；Muon 用 NS 高效近似极分解正交因子，提供与 Adam 不同的\textbf{非对角}预条件 $P_k$。
  \item \textbf{本报告实验目的}：验证上述几何解释（NS 收敛盆、trust-region 方向、Frobenius 错配），而非在简化二次模型上宣称 Muon 优于 GD。
\end{enumerate}
